\documentclass{article}
\usepackage[svgnames]{xcolor}
\usepackage{bm}
\usepackage{graphicx}
\usepackage{amsfonts, amsmath, amsthm, amssymb}
\usepackage{breqn}
\usepackage{empheq}
\usepackage{listings}
\usepackage{booktabs}
\usepackage{makecell}
\usepackage{float}
\usepackage{longtable}
\usepackage{geometry}
\usepackage{enumitem}
\theoremstyle{plain}
\usepackage{tcolorbox}
\usepackage{pgfplots}
\usepackage{multicol}
\pgfplotsset{compat=1.18}
\usepackage{siunitx}
\usepackage{listings}
\usepackage{tabularx}
\usepackage{rotating}
\usepackage{longtable}
\usepackage{mathtools}
\usepackage{geometry}  % For better page margins (optional)
\geometry{a4paper, margin=1in} 
\newtheorem{proposition}{Proposition}[section]
\theoremstyle{remark}
\newtheorem*{remark}{Remark}  % Asterisk prevents numbered remarks
\usepackage{tocloft} % Required for \cft commands
\usepackage{xcolor} % Required for colors
\usepackage[numbers]{natbib}
\bibliographystyle{plainnat} % or abbrvnat, unsrtnat, etc.
\usepackage{titlesec} % For customizing section titles
% Define a new 'subsubsubsection' level (level 4)
\titleclass{\subsubsubsection}{straight}[\subsection]
\newcounter{subsubsubsection}[subsubsection]
\renewcommand{\thesubsubsubsection}{\thesubsubsection.\arabic{subsubsubsection}}
\titleformat{\subsubsubsection}{\normalfont\normalsize\bfseries}{\thesubsubsubsection}{1em}{}
\titlespacing*{\subsubsubsection}{0pt}{3.25ex plus 1ex minus .2ex}{1.5ex plus .2ex}
\newtheorem{example}{Example} % Numbered examples
% Define a custom theorem style for conjectures
\newtheoremstyle{myconjecture}
  {10pt}   % Space above
  {10pt}   % Space below
  {\itshape}  % Body font
  {0pt}    % Indent
  {\bfseries} % Head font
  {.}      % Punctuation after head
  { }      % Space after head
  {\thmname{#1}\thmnumber{ #2}\thmnote{ (#3)}} % Custom head spec

\theoremstyle{myconjecture}
\newtheorem{conjecture}{Conjecture}

% Set numbering depth
\setcounter{secnumdepth}{4} % Number down to subsubsubsection
\setcounter{tocdepth}{4}    % Show down to subsubsubsection in TOC



% Theorem Environments
\theoremstyle{plain}
\newtheorem{exob}{Experimental Observation}
\newtheorem{theorem}{Theorem}[section]
\newtheorem{lemma}[theorem]{Lemma}
\newtheorem{corollary}[theorem]{Corollary}
\theoremstyle{definition}
\newenvironment{keyword}
{\noindent\textbf{Keywords:}\ }
{\par\medskip}

% TOC Customization
\renewcommand{\cftsecfont}{\normalfont\bfseries}
\renewcommand{\cftsecleader}{\cftdotfill{\cftdotsep}}
\setlength{\cftbeforesecskip}{4pt}


% Hyperref (LAST package)
\usepackage{hyperref}
\hypersetup{colorlinks=true, linkcolor=blue}
\usepackage{cleveref} % Auto-reference "Equation (X)"
\numberwithin{equation}{section} % Reset numbering per section
\usepackage{tikz}
\usetikzlibrary{arrows.meta, positioning}

% Define colors and style
\definecolor{codegreen}{rgb}{0,0.6,0}
\definecolor{codegray}{rgb}{0.5,0.5,0.5}
\definecolor{codepurple}{rgb}{0.58,0,0.82}
\definecolor{backcolour}{rgb}{0.95,0.95,0.92}

\lstdefinestyle{sagemath}{
    backgroundcolor=\color{backcolour},   
    commentstyle=\color{codegreen},
    keywordstyle=\color{magenta},
    numberstyle=\tiny\color{codegray},
    stringstyle=\color{codepurple},
    basicstyle=\ttfamily\footnotesize,
    breakatwhitespace=false,         
    breaklines=true,                 
    captionpos=b,                    
    keepspaces=true,                 
    numbers=left,                    
    numbersep=5pt,                  
    showspaces=false,                
    showstringspaces=false,
    showtabs=false,                  
    tabsize=2,
    frame=single,
    rulecolor=\color{black},
    language=Python,
    morekeywords={True,False,sage,for,while,if,elif,else,def,return,print,lambda,in,try,except,with,as,assert,class,from,import,pass,break,continue,global,nonlocal,None,not,and,or,is},
}


\title{Infinitely Many Integer Solutions for Five Related Diophantine Equations}
\author{Jamal Agbanwa \and Bogdan Grechuk}
\date{}

\begin{document}
\maketitle

\begin{abstract}
We prove that each of the five equations
\[
 y^2+x^3y+z^2-2=0,\qquad
 y^2+x^3y+z^2+z-1=0,
\]
\[
 y^2+x^3y+z^2+z+1=0,\qquad
 y^2+x^3y+y+z^2+1=0,
\]
\[
 y^2+x^3y+z^2+1=0
\]
has infinitely many integer solutions.  The proof is constructive: explicit
Pell-type recurrences produce infinitely many distinct values of the parameter
\(x\).  A tangent identity on the surface of sums of two squares supplies
rational two-square representations, and the classical two-squares theorem
converts those rational representations into integral ones.
\end{abstract}

\noindent\textbf{Scope of the result.}
The paper proves infinitude of integer solutions for the five displayed
equations.  It does not claim to classify all integer solutions.  Each
construction below gives an infinite subfamily of solutions, distinguished by
infinitely many different values of \(x\).


\noindent\textbf{Background.}
The five equations treated here are taken from Table~13 of Grechuk's list of
short open Diophantine problems, where they appear among the equations of
length \(\ell\le 9\) for which Problem~4 was open.  In that reference,
Problem~4 asks one either to list all integer solutions of a given Diophantine
equation or to prove that it has infinitely many integer solutions
\cite[Problem~4 and Table~13]{Grechuk2024}.

\section{Two elementary tools}

We shall use the classical two-squares theorem.  A proof is included for
self-containment.

\begin{theorem}[Fermat's two-squares theorem]\label{thm:two-squares}
A positive integer \(m\) is a sum of two integer squares if and only if every
prime \(p\equiv3\pmod4\) occurs to an even exponent in the prime factorization
of \(m\).
\end{theorem}

\begin{proof}
We use the Gaussian integers \(\mathbb Z[i]\).  For
\(\alpha=r+si\in\mathbb Z[i]\), put
\[
        N(\alpha)=\alpha\overline\alpha=r^2+s^2.
\]
This norm is multiplicative, and its units are exactly the elements of norm
\(1\), namely \(\pm1,\pm i\).

First recall that \(\mathbb Z[i]\) is Euclidean.  If
\(\alpha,\beta\in\mathbb Z[i]\) and \(\beta\ne0\), choose a Gaussian integer
\(q\) whose real and imaginary parts are nearest integers to those of
\(\alpha/\beta\).  Put \(r=\alpha-q\beta\).  Then
\[
        N(r)=N(\beta)\left|\frac{\alpha}{\beta}-q\right|^2
        \le \frac{1}{2}N(\beta)<N(\beta).
\]
Thus Euclidean division is available.  Hence every ideal of \(\mathbb Z[i]\)
is principal: if \(I\ne0\) is an ideal, choose \(\delta\in I\setminus\{0\}\)
with minimal norm.  For any \(\alpha\in I\), write
\(\alpha=q\delta+r\), with either \(r=0\) or \(N(r)<N(\delta)\).  Since
\(r\in I\), minimality forces \(r=0\), so \(I=(\delta)\).  Therefore
\(\mathbb Z[i]\) is a principal ideal domain.

In a principal ideal domain, every irreducible element is prime.  Indeed, if
\(\pi\) is irreducible and \(\pi\mid\alpha\beta\) but \(\pi\nmid\alpha\), then
\((\pi,\alpha)=(1)\), so there exist \(r,s\in\mathbb Z[i]\) with
\(r\pi+s\alpha=1\).  Multiplying by \(\beta\) shows \(\pi\mid\beta\).
Existence of factorizations follows by norm descent: if a nonzero nonunit is
not irreducible, write it as a product of two nonunits, each of smaller
positive norm, and repeat.  The process terminates because positive integers
cannot decrease forever.  Uniqueness follows by induction on the number of
irreducible factors, using the prime property to cancel one irreducible factor
at a time.  Thus \(\mathbb Z[i]\) is a unique factorization domain.

We next classify rational primes in \(\mathbb Z[i]\).  The prime \(2\) ramifies:
\[
        2=(1+i)(1-i).
\]
Let \(p\) be an odd rational prime.  If \(p\equiv3\pmod4\), then \(-1\) is not a
quadratic residue modulo \(p\).  For if \(t^2\equiv-1\pmod p\), Fermat's theorem
would give
\[
        1\equiv t^{p-1}=(t^2)^{(p-1)/2}
        \equiv (-1)^{(p-1)/2}\equiv -1\pmod p,
\]
because \((p-1)/2\) is odd.  Here Fermat's theorem follows from the fact that
multiplication by \(t\) permutes the nonzero residue classes modulo \(p\).
Thus \(X^2+1\) is irreducible over \(\mathbb F_p\), and
\[
        \mathbb Z[i]/(p)\cong \mathbb F_p[X]/(X^2+1)
\]
is a field.  Hence \((p)\) is a prime ideal, equivalently \(p\) is prime in
\(\mathbb Z[i]\).

If \(p\equiv1\pmod4\), then \(-1\) is a quadratic residue modulo \(p\).  Wilson's
theorem for odd primes gives \((p-1)!\equiv-1\pmod p\).  Also
\[
        (p-1)!\equiv
        \prod_{j=1}^{(p-1)/2} j(p-j)
        \equiv (-1)^{(p-1)/2}\left(\left(\frac{p-1}{2}\right)!\right)^2
        \pmod p.
\]
Since \(p\equiv1\pmod4\), the exponent \((p-1)/2\) is even, and hence
\[
        \left(\left(\frac{p-1}{2}\right)!\right)^2\equiv -1\pmod p.
\]
Wilson's theorem itself follows by multiplying all nonzero residue classes
modulo \(p\): every class pairs with its inverse, except \(1\) and \(-1\), so
the product is \(-1\).  Therefore choose \(t\) with \(t^2\equiv-1\pmod p\).
Then \(p\mid(t+i)(t-i)\) in \(\mathbb Z[i]\), but \(p\) divides neither factor.
Thus \(p\) is not prime in \(\mathbb Z[i]\).  Since irreducibles are prime in a
UFD, \(p\) is not irreducible.  Hence \(p=\alpha\beta\) for nonunits
\(\alpha,\beta\in\mathbb Z[i]\).  Taking norms gives
\[
        p^2=N(p)=N(\alpha)N(\beta).
\]
The two factors on the right are integers greater than \(1\).  Since \(p\) is a
rational prime, the only possibility is \(N(\alpha)=N(\beta)=p\).  Put
\(\pi=\alpha\).  Then \(N(\pi)=p\), and therefore
\[
        p=\pi\overline\pi.
\]
Thus a rational prime \(p\equiv1\pmod4\) splits into two conjugate Gaussian
factors, each of norm \(p\).

Suppose first that \(m=r^2+s^2\).  Then
\[
        m=(r+si)(r-si).
\]
Let \(p\equiv3\pmod4\).  As just shown, \(p\) is prime in \(\mathbb Z[i]\), and
it is associated to its conjugate.  Therefore its exponent in the product
\((r+si)(r-si)\) is twice its exponent in \(r+si\), hence is even.  Thus every
prime \(p\equiv3\pmod4\) occurs to an even exponent in \(m\).

Conversely, suppose every prime \(p\equiv3\pmod4\) occurs to an even exponent in
\(m\).  Write the prime factorization of \(m\) in \(\mathbb Z\).  For each prime
\(p\equiv1\pmod4\), choose \(\pi_p\in\mathbb Z[i]\) with \(N(\pi_p)=p\).  For
\(p=2\), use \(N(1+i)=2\).  For each prime \(q\equiv3\pmod4\), the exponent of
\(q\) in \(m\) is even, say \(2e_q\), and \(N(q^{e_q})=q^{2e_q}\).  Hence the
Gaussian integer
\[
        \alpha=(1+i)^{v_2(m)}
        \prod_{p\equiv1\,({\rm mod}\,4)}\pi_p^{v_p(m)}
        \prod_{q\equiv3\,({\rm mod}\,4)}q^{v_q(m)/2}
\]
has norm \(m\).  Writing \(\alpha=r+si\), we obtain
\[
        m=N(\alpha)=r^2+s^2.
\]
This proves the theorem.
\end{proof}

\begin{lemma}[Rational-to-integral two-square lemma]\label{lem:rational-integral}
Let \(m\) be a positive integer.  If
\[
        m=A^2+B^2
\]
for some \(A,B\in\mathbb Q\), then
\[
        m=r^2+s^2
\]
for some \(r,s\in\mathbb Z\).
\end{lemma}

\begin{proof}
Choose a positive integer \(q\) such that \(qA,qB\in\mathbb Z\).  Then
\[
        mq^2=(qA)^2+(qB)^2
\]
is a sum of two integer squares.  By Theorem \ref{thm:two-squares}, every prime
\(p\equiv3\pmod4\) occurs to an even exponent in \(mq^2\).  Since \(q^2\)
contributes only even exponents, every prime \(p\equiv3\pmod4\) already occurs
to an even exponent in \(m\).  Applying Theorem \ref{thm:two-squares} again,
\(m\) is a sum of two integer squares.
\end{proof}

We shall also use the Brahmagupta-Fibonacci identity
\[
        (r^2+s^2)(p^2+q^2)=(rp-sq)^2+(rq+sp)^2.       \tag{1}\label{eq:bf}
\]

\section{A tangent identity for \texorpdfstring{\(N^6+c\)}{N6+c}}

The next lemma is the main source of sums of two squares.

\begin{lemma}[Tangent identity]\label{lem:tangent}
Let \(c,u,a,b\in\mathbb Z\), and put
\[
        S=u^3+c.
\]
Assume that \(S\ne0\) and
\[
        a^2+b^2=S.
\]
If \(N,R\in\mathbb Z\) satisfy
\[
        R^2=4SN^2-u(u^3-8c),                 \tag{2}\label{eq:Rcond}
\]
then \(N^6+c\) is a sum of two rational squares.
\end{lemma}

\begin{proof}
The hypotheses imply \(S=a^2+b^2>0\), so the divisions by \(2S\) below are
legitimate.  Set
\[
        L=3u^2N^2-u^3+2c.
\]
We first verify the polynomial identity
\[
        4S(N^6+c)-L^2
        =(N^2-u)^2\{4SN^2-u(u^3-8c)\}.            \tag{3}\label{eq:tangent-id}
\]
Put \(T=N^2\).  The left-hand side of \eqref{eq:tangent-id} equals
\[
\begin{aligned}
&4(u^3+c)(T^3+c)-\bigl(3u^2T-u^3+2c\bigr)^2  \\
&\quad=4(u^3+c)T^3-9u^4T^2+(6u^5-12u^2c)T-u^6+8cu^3.
\end{aligned}
\]
The right-hand side equals
\[
\begin{aligned}
&(T-u)^2\{4(u^3+c)T-u(u^3-8c)\} \\
&\quad=(T^2-2uT+u^2)(4(u^3+c)T-u^4+8uc) \\
&\quad=4(u^3+c)T^3-9u^4T^2+(6u^5-12u^2c)T-u^6+8cu^3.
\end{aligned}
\]
Thus \eqref{eq:tangent-id} is proved.

Using \eqref{eq:Rcond}, identity \eqref{eq:tangent-id} becomes
\[
        4S(N^6+c)-L^2=(N^2-u)^2R^2.
\]
Define rational numbers
\[
        A=\frac{aL-b(N^2-u)R}{2S},
        \qquad
        B=\frac{bL+a(N^2-u)R}{2S}.
\]
Since \(S=a^2+b^2\), the mixed terms cancel and give
\[
\begin{aligned}
        A^2+B^2
        &=\frac{L^2+(N^2-u)^2R^2}{4S} \\
        &=\frac{4S(N^6+c)}{4S}
        =N^6+c.
\end{aligned}
\]
Thus \(N^6+c\) is a sum of two rational squares.
\end{proof}

The Pell recurrences used below are packaged in the following elementary form.

\begin{lemma}[Pell propagation]\label{lem:pell}
Let \(D\) be a positive nonsquare integer, and let \(K\in\mathbb Z\).  Suppose
\(X_0,N_0,P,Q\) are positive integers satisfying
\[
        X_0^2-DN_0^2=K,
        \qquad
        P^2-DQ^2=1.
\]
Define
\[
        X_{k+1}=PX_k+DQN_k,
        \qquad
        N_{k+1}=QX_k+PN_k
        \qquad(k\ge0).
\]
Then
\[
        X_k^2-DN_k^2=K
\]
for every \(k\ge0\).  Moreover, the positive integers \(N_k\) are pairwise
distinct.  If \(P\) is odd and \(Q\) is even, then the parities of \(X_k\) and
\(N_k\) are independent of \(k\).
\end{lemma}

\begin{proof}
Equivalently,
\[
        X_k+N_k\sqrt D=(X_0+N_0\sqrt D)(P+Q\sqrt D)^k.
\]
Taking norms gives \(X_k^2-DN_k^2=K\).  Since \(P^2=1+DQ^2>1\), we have
\(P>1\).  Positivity is clear by induction, and
\[
        N_{k+1}=QX_k+PN_k>N_k.
\]
Thus the positive integers \(N_k\) are pairwise distinct.  Finally, if \(P\) is
odd and \(Q\) is even, then
\[
        X_{k+1}=PX_k+DQN_k\equiv X_k\pmod2,
        \qquad
        N_{k+1}=QX_k+PN_k\equiv N_k\pmod2.
\]
\end{proof}

\section{Four infinite families of two-square values}

\begin{proposition}\label{prop:c8}
There are infinitely many positive even integers \(N\) for which \(N^6+8\) is a
sum of two integer squares.
\end{proposition}

\begin{proof}
Apply Lemma \ref{lem:tangent} with
\[
        c=8,
        \qquad u=8,
        \qquad a=6,
        \qquad b=22.
\]
Indeed,
\[
        6^2+22^2=520=8^3+8.
\]
For these constants, condition \eqref{eq:Rcond} is
\[
        R^2=2080N^2-3584=16(130N^2-224).
\]
Thus it is enough to solve
\[
        X^2-130N^2=-224,                    \tag{4}\label{eq:pell-c8}
\]
and then put \(R=4X\).  The initial solution and a Pell unit are
\[
        136^2-130\cdot 12^2=-224,
        \qquad
        6499^2-130\cdot570^2=1.
\]
Lemma \ref{lem:pell}, with
\[
        D=130,
        \quad K=-224,
        \quad X_0=136,
        \quad N_0=12,
        \quad P=6499,
        \quad Q=570,
\]
produces infinitely many positive integers \(N_k\) solving
\eqref{eq:pell-c8}.  Since \(N_0\) is even, \(P\) is odd, and \(Q\) is even,
the parity assertion in Lemma \ref{lem:pell} shows that every \(N_k\) is even.
Lemma \ref{lem:tangent} gives rational representations of \(N_k^6+8\) as sums
of two squares, and Lemma \ref{lem:rational-integral} converts them to integer
representations.
\end{proof}

\begin{proposition}\label{prop:c5}
There are infinitely many positive even integers \(N\) for which \(N^6+5\) is a
sum of two integer squares.
\end{proposition}

\begin{proof}
Apply Lemma \ref{lem:tangent} with
\[
        c=5,
        \qquad u=2,
        \qquad a=2,
        \qquad b=3.
\]
Here
\[
        2^2+3^2=13=2^3+5.
\]
The condition \eqref{eq:Rcond} becomes
\[
        R^2=52N^2+64=4(13N^2+16).
\]
Thus it is enough to solve
\[
        X^2-13N^2=16,                       \tag{5}\label{eq:pell-c5}
\]
and put \(R=2X\).  We have
\[
        22^2-13\cdot6^2=16,
        \qquad
        649^2-13\cdot180^2=1.
\]
Lemma \ref{lem:pell}, with
\[
        D=13,
        \quad K=16,
        \quad X_0=22,
        \quad N_0=6,
        \quad P=649,
        \quad Q=180,
\]
therefore gives infinitely many positive integers \(N_k\) solving
\eqref{eq:pell-c5}.  Since \(N_0\) is even, \(P\) is odd, and \(Q\) is even,
the same lemma shows that every \(N_k\) is even.  Lemma \ref{lem:tangent} and
Lemma \ref{lem:rational-integral} then imply that each \(N_k^6+5\) is a sum of
two integer squares.
\end{proof}

\begin{proposition}\label{prop:cm3}
There are infinitely many positive even integers \(N\) for which \(N^6-3\) is a
sum of two integer squares.
\end{proposition}

\begin{proof}
Apply Lemma \ref{lem:tangent} with
\[
        c=-3,
        \qquad u=2,
        \qquad a=1,
        \qquad b=2.
\]
Indeed,
\[
        1^2+2^2=5=2^3-3.
\]
The condition \eqref{eq:Rcond} becomes
\[
        R^2=20N^2-64=4(5N^2-16).
\]
Thus it is enough to solve
\[
        X^2-5N^2=-16,                       \tag{6}\label{eq:pell-cm3}
\]
and put \(R=2X\).  We have
\[
        2^2-5\cdot2^2=-16,
        \qquad
        9^2-5\cdot4^2=1.
\]
Lemma \ref{lem:pell}, with
\[
        D=5,
        \quad K=-16,
        \quad X_0=2,
        \quad N_0=2,
        \quad P=9,
        \quad Q=4,
\]
gives infinitely many positive integers \(N_k\) solving
\eqref{eq:pell-cm3}.  Since \(N_0\) is even, \(P\) is odd, and \(Q\) is even,
Lemma \ref{lem:pell} shows that every \(N_k\) is even.  Lemma
\ref{lem:tangent} and Lemma \ref{lem:rational-integral} then imply that each
\(N_k^6-3\) is a sum of two integer squares.
\end{proof}

\begin{proposition}\label{prop:cm4}
There are infinitely many positive odd integers \(N\) for which \(N^6-4\) is a
sum of two integer squares.
\end{proposition}

\begin{proof}
Apply Lemma \ref{lem:tangent} with
\[
        c=-4,
        \qquad u=242,
        \qquad a=578,
        \qquad b=3720.
\]
Then
\[
        a^2+b^2=578^2+3720^2=14172484=242^3-4.
\]
Thus, in the notation of Lemma \ref{lem:tangent},
\(S=14172484\ne0\).  The condition \eqref{eq:Rcond} becomes
\[
        R^2=56689936N^2-3429749840.
\]
Since
\[
        56689936=16\cdot3543121,
        \qquad
        3429749840=16\cdot214359365,
\]
it is enough to solve
\[
        X^2-3543121N^2=-214359365,          \tag{7}\label{eq:pell-cm4}
\]
and then put \(R=4X\).  The integer \(3543121\) is nonsquare, since
\[
        1882^2<3543121<1883^2.
\]
The following two identities are direct integer equalities:
\[
\begin{aligned}
&410154078658987088294^2
 -3543121\cdot217898400660155259^2=-214359365, \\
&53319160674725436041296699414748449^2 \\
&\qquad
 -3543121\cdot28326330128303226243386806483560^2=1.
\end{aligned}
\]
Now apply Lemma \ref{lem:pell} with
\[
\begin{gathered}
        D=3543121,\qquad K=-214359365, \\
        X_0=410154078658987088294,\qquad
        N_0=217898400660155259, \\
        P=53319160674725436041296699414748449, \\
        Q=28326330128303226243386806483560.
\end{gathered}
\]
This produces infinitely many positive integers \(N_k\) solving
\eqref{eq:pell-cm4}.  The parity is preserved because \(P\) is odd and \(Q\) is
even.  Since \(N_0\) is odd, every \(N_k\) is odd.  Lemma \ref{lem:tangent}
gives rational representations of \(N_k^6-4\) as sums of two squares, and Lemma
\ref{lem:rational-integral} converts them to integer representations.
\end{proof}

\section{Converting two-square values into solutions}

We now turn the two-square values above into solutions of the five equations.

\begin{theorem}\label{thm:first-three}
Each of the equations
\[
        y^2+x^3y+z^2-2=0,                    \tag{E1}\label{eq:E1}
\]
\[
        y^2+x^3y+z^2+z-1=0,                  \tag{E2}\label{eq:E2}
\]
\[
        y^2+x^3y+z^2+z+1=0                   \tag{E3}\label{eq:E3}
\]
has infinitely many integer solutions.
\end{theorem}

\begin{proof}
For \eqref{eq:E1}, let \(N\) run through the infinitely many positive even
integers supplied by Proposition \ref{prop:c8}.  For each such \(N\), choose
integers \(a,b\) with
\[
        N^6+8=a^2+b^2.
\]
Since \(N\) is even, \(N^6+8\equiv8\pmod {16}\).  A square modulo \(16\) is one
of \(0,1,4,9\), so a representation of a number congruent to \(8\pmod {16}\)
forces both \(a\) and \(b\) to be even.  Put
\[
        x=N,
        \qquad
        y=\frac{a-N^3}{2},
        \qquad
        z=\frac b2.
\]
Then \(y,z\in\mathbb Z\), and
\[
        (2y+x^3)^2+(2z)^2=a^2+b^2=N^6+8=x^6+8.
\]
Expanding and dividing by \(4\) gives
\[
        y^2+x^3y+z^2-2=0.
\]
The values \(x=N\) are pairwise distinct, so \eqref{eq:E1} has infinitely many
integer solutions.

For \eqref{eq:E2}, let \(N\) run through the infinitely many positive even
integers supplied by Proposition \ref{prop:c5}.  Choose integers \(a,b\) with
\[
        N^6+5=a^2+b^2.
\]
Since \(N\) is even, \(N^6+5\equiv5\pmod8\).  Hence one of \(a,b\) is even and
the other is odd.  Relabel them, if necessary, so that \(a\) is even and \(b\)
is odd.  Put
\[
        x=N,
        \qquad
        y=\frac{a-N^3}{2},
        \qquad
        z=\frac{b-1}{2}.
\]
Then \(y,z\in\mathbb Z\), and
\[
        (2y+x^3)^2+(2z+1)^2=a^2+b^2=N^6+5=x^6+5.
\]
Expanding, cancelling \(x^6\), and dividing by \(4\) gives
\[
        y^2+x^3y+z^2+z-1=0.
\]
Again the values \(x=N\) are pairwise distinct, so \eqref{eq:E2} has infinitely
many integer solutions.

For \eqref{eq:E3}, use the infinitely many positive even integers \(N\) supplied
by Proposition \ref{prop:cm3}.  Choose integers \(a,b\) with
\[
        N^6-3=a^2+b^2.
\]
Since \(N\) is even, \(N^6-3\equiv5\pmod8\).  Thus one of \(a,b\) is even and
the other is odd; relabel so that \(a\) is even and \(b\) is odd.  Put
\[
        x=N,
        \qquad
        y=\frac{a-N^3}{2},
        \qquad
        z=\frac{b-1}{2}.
\]
Then \(y,z\in\mathbb Z\), and
\[
        (2y+x^3)^2+(2z+1)^2=a^2+b^2=N^6-3=x^6-3,
\]
and expanding, cancelling \(x^6\), and dividing by \(4\) gives
\[
        y^2+x^3y+z^2+z+1=0.
\]
The values \(x=N\) are pairwise distinct, so \eqref{eq:E3} has infinitely many
integer solutions.
\end{proof}

\begin{theorem}\label{thm:fourth}
The equation
\[
        y^2+x^3y+y+z^2+1=0                  \tag{E4}\label{eq:E4}
\]
has infinitely many integer solutions.
\end{theorem}

\begin{proof}
Let \(N\) run through the infinitely many positive even integers supplied by
Proposition \ref{prop:cm3}.  Thus
\[
        N^6-3=r^2+s^2
\]
for some integers \(r,s\).  Also
\[
        N^6+1=(N^3)^2+1^2.
\]
By \eqref{eq:bf}, the product
\[
        (N^6-3)(N^6+1)
\]
is a sum of two integer squares.  Now set
\[
        x=-N^2.
\]
Then
\[
\begin{aligned}
        (x^3+1)^2-4
        &=(1-N^6)^2-4 \\
        &=N^{12}-2N^6-3 \\
        &=(N^6-3)(N^6+1).
\end{aligned}
\]
Therefore there exist integers \(a,b\) such that
\[
        (x^3+1)^2-4=a^2+b^2.
\]
Because \(N\) is even, \(x=-N^2\) is even, so \(x^3+1\) is odd and
\[
        (x^3+1)^2-4\equiv5\pmod8.
\]
Thus one of \(a,b\) is odd and the other is even.  Relabel, if necessary, so
that \(a\) is odd and \(b\) is even.  Put
\[
        y=\frac{a-(x^3+1)}{2},
        \qquad
        z=\frac b2.
\]
Then \(y,z\in\mathbb Z\), and
\[
        (2y+x^3+1)^2+(2z)^2=a^2+b^2=(x^3+1)^2-4.
\]
Expanding, cancelling \((x^3+1)^2\), and dividing by \(4\) gives
\[
        y^2+x^3y+y+z^2+1=0.
\]
The values \(x=-N^2\) are pairwise distinct as \(N\) varies, so \eqref{eq:E4}
has infinitely many integer solutions.
\end{proof}

\begin{theorem}\label{thm:fifth}
The equation
\[
        y^2+x^3y+z^2+1=0                    \tag{E5}\label{eq:E5}
\]
has infinitely many integer solutions.
\end{theorem}

\begin{proof}
Let \(N\) run through the infinitely many positive odd integers supplied by
Proposition \ref{prop:cm4}.  For each such \(N\), choose integers \(a,b\) with
\[
        N^6-4=a^2+b^2.
\]
Since \(N\) is odd,
\[
        N^6-4\equiv 1-4\equiv5\pmod8.
\]
The quadratic residues modulo \(8\) are \(0,1,4\).  Therefore one of \(a,b\) is
odd and the other is even.  Relabel them, if necessary, so that \(a\) is odd
and \(b\) is even.  Put
\[
        x=N,
        \qquad
        y=\frac{a-N^3}{2},
        \qquad
        z=\frac b2.
\]
These are integers, because \(a\) and \(N^3\) are odd and \(b\) is even.  Also
\[
        (2y+x^3)^2+(2z)^2=a^2+b^2=N^6-4=x^6-4.
\]
Expanding the left-hand side gives
\[
        4y^2+4x^3y+x^6+4z^2=x^6-4.
\]
After cancellation and division by \(4\), this is
\[
        y^2+x^3y+z^2+1=0.
\]
The values \(x=N\) are pairwise distinct, so \eqref{eq:E5} has infinitely many
integer solutions.
\end{proof}

Combining Theorems \ref{thm:first-three}, \ref{thm:fourth}, and
\ref{thm:fifth} proves the desired claim for all five equations.

\begin{remark}
The construction gives explicit infinite families.  The \(N\)'s in Propositions
\ref{prop:c8}, \ref{prop:c5}, \ref{prop:cm3}, and \ref{prop:cm4} are obtained
from the stated Pell recurrences.  The proof does not claim to classify all
integer solutions; it only supplies infinitely many of them.
\end{remark}

\appendix
\section{Exact arithmetic certificate for the large Pell identities}

The proof of Proposition \ref{prop:cm4} uses two large Pell-type identities.
They are ordinary integer equalities, not heuristic or floating-point checks.
The following plain Python 3 script verifies them using exact arbitrary-precision
integer arithmetic.

\begin{lstlisting}[caption={Exact verification of the large Pell norm identities.}]
D = 3543121
C = 214359365

X0 = 410154078658987088294
N0 = 217898400660155259

P = 53319160674725436041296699414748449
Q = 28326330128303226243386806483560

initial_norm = X0*X0 - D*N0*N0
unit_norm = P*P - D*Q*Q

assert initial_norm == -C
assert unit_norm == 1

# Parity data used in the c = -4 construction.
assert X0 % 2 == 0
assert N0 % 2 == 1
assert P % 2 == 1
assert Q % 2 == 0

print("initial_norm =", initial_norm)
print("unit_norm =", unit_norm)
print("parities (X0,N0,P,Q) =", (X0 % 2, N0 % 2, P % 2, Q % 2))
\end{lstlisting}

The expected output is
\begin{lstlisting}
initial_norm = -214359365
unit_norm = 1
parities (X0,N0,P,Q) = (0, 1, 1, 0)
\end{lstlisting}

For completeness, the smaller constants used to specialize Lemma
\ref{lem:tangent} in Proposition \ref{prop:cm4} may be checked in the same way:

\begin{lstlisting}[caption={Exact verification of the constants in Proposition \ref{prop:cm4}.}]
u = 242
S = u**3 - 4

assert S == 14172484
assert 578**2 + 3720**2 == S
assert 4*S == 16*3543121
assert u*(u**3 + 32) == 16*214359365

print("small constants verified")
\end{lstlisting}

These scripts are included only as checkable certificates for arithmetic
substitutions.  The mathematical proof of infinitude is the Pell recurrence and
norm argument given in Lemma \ref{lem:pell} and Proposition \ref{prop:cm4}.


\bibliography{references}

\vspace{2em}

\noindent\textbf{Jamal Agbanwa}\\
Independent Researcher, Rotterdam, The Netherlands\\
\href{mailto:agbanwajamal03@gmail.com}{\texttt{agbanwajamal03@gmail.com}}

\vspace{1em}

\noindent\textbf{Bogdan Grechuk}\\
Department of Mathematics, University of Leicester, Leicester, LE1 7RH, UK\\
\href{mailto:bg83@leicester.ac.uk}{\texttt{bg83@leicester.ac.uk}}

\end{document}
